\section{Thiết lập tái lập}\label{sec:setup}

Mục này mô tả các điều kiện dùng chung cho mọi phép đo trong bài.
Từ đây trở đi, mọi cấu hình chỉ khác nhau ở hàm xếp hạng đoạn (\cref{sec:design}).

\subsection{Dữ liệu}
LiHua-World~\cite{fan2025minirag} gồm 442 tài liệu hội thoại tiếng Anh của một người dùng giả định. Mỗi tài liệu mang mốc thời gian.
Bộ câu hỏi có 637 dòng (635 câu phân biệt) thuộc ba loại: \emph{Single} (một bước), \emph{Multi} (nhiều bước) và \emph{Null} (đáp án đúng là ``không đủ thông tin'').
Cột \emph{Evidence} ghi mốc thời gian của các tin nhắn chứa đáp án, nhờ đó ta xác định được chính xác đoạn nào chứa bằng chứng.
Chúng tôi chỉ dùng cột này để chấm truy xuất; không hệ thống nào đọc nó lúc chạy.

Chúng tôi tách 200 câu làm \emph{tập dev}, phân tầng theo loại với seed cố định.
Mọi chẩn đoán, sàng lọc và lựa chọn cấu hình đều làm trên tập dev.
Còn lại \nHeld{} câu (\nHeldSingle{} Single, \nHeldMulti{} Multi, \nHeldNull{} Null) tạo thành \emph{tập ngoài dev}, chỉ dùng cho vòng xác nhận.

\subsection{Hệ thống}
Mô hình \texttt{Qwen/Qwen2.5-3B-Instruct}~\cite{qwen2025qwen25} làm mọi việc cần mô hình ngôn ngữ: trích xuất thực thể khi lập chỉ mục, phân tích câu hỏi và sinh câu trả lời.
Mô hình được phục vụ bằng vLLM 0.11.0 với cửa sổ 32.768 token; nhiệt độ lấy mẫu không đặt trong yêu cầu nên dùng giá trị mặc định của máy chủ \todo{xác nhận nhiệt độ thực tế}.
Vectơ nhúng là \texttt{all-MiniLM-L6-v2}~\cite{reimers2019sbert,wang2020minilm} (384 chiều), giống bài gốc.
Đoạn văn bản dài tối đa 1.200 token. Đồ thị do chính Qwen2.5-3B dựng có \graphNodes{} nút và \graphEdges{} cạnh, trong đó \graphIsolated\% nút không có cạnh nào.
Tham số truy vấn là \texttt{top\_k} = 60 và ngân sách $\budget = 4000$ token cho phần đoạn văn bản.
Mọi cấu hình dùng cùng một chỉ mục đã đóng băng.

\subsection{Hai lỗi phải sửa để chạy với mô hình nhỏ}\label{sec:fixes}
Mã gốc có hai lỗi khiến nó không chạy được với SLM trong ngân sách hợp lý.
Chúng tôi coi việc sửa là \emph{điều kiện để tái lập}, không phải đóng góp, và áp dụng như nhau cho mọi cấu hình.
\begin{itemize}
  \item \textbf{Lập chỉ mục lặp $O(N^2)$.} Hàm chèn tài liệu dựng lại tập đoạn từ mọi tài liệu đã xử lý, nên chèn $N$ tài liệu tuần tự sẽ trích xuất lại toàn bộ kho $N$ lần.
  Lưu vết các đoạn đã trích xuất đưa số lời gọi mô hình từ khoảng \callsBefore{} xuống \callsAfter{} cho 267 tài liệu.
  \item \textbf{Ngân sách ngữ cảnh không giới hạn.} Trên 200 câu dev, phần đoạn của ngữ cảnh có trung vị \ctxNoCapMed{} token và tối đa \ctxNoCapMax{} token, vượt cửa sổ của Qwen2.5-3B ở một số câu.
  Chúng tôi cắt phần đoạn ở $\budget = 4000$ token, đưa trung vị xuống \ctxCapMed{} token.
\end{itemize}
Cơ chế khớp loại đáp án của MiniRAG cũng không khớp được nút nào vì so sánh phân biệt hoa thường; sau khi sửa, nó khớp \atypeNodesFixed{} nút.
Bản sửa này không làm accuracy thay đổi phân biệt được (Qwen: \atypeQwen{}; Gemini: \atypeGemini{}; McNemar, 635 câu), nhưng chúng tôi vẫn giữ nó vì nó sửa một lỗi thật.

\subsection{Đánh giá}\label{sec:eval}
Giám khảo \texttt{gemini-flash-lite-latest} gán cho mỗi câu trả lời một trong ba nhãn, theo giao thức của bài gốc: \emph{accurate} (truyền đạt đúng đáp án vàng), \emph{error} (khẳng định điều mâu thuẫn với đáp án vàng mà không rào đón) và \emph{neither} (từ chối, không biết hoặc lạc đề).
Prompt đầy đủ ở \cref{app:prompts}.
Chúng tôi báo tỉ lệ của cả ba nhãn và tỉ số acc/(acc+err); tỉ số này tách hiệu ứng ``dám trả lời nhiều hơn'' khỏi hiệu ứng ``trả lời đúng hơn''.
Với câu Null, lời từ chối là câu trả lời đúng.

Nhiễu giữa các lượt sinh lớn hơn nhiều so với nhiễu giám khảo.
Trên MiniRAG tái lập, độ lệch chuẩn accuracy giữa ba lượt chấm của cùng một lượt sinh là \sdJudgeBase{} điểm; giữa ba lượt sinh của RRF(đồ thị, vectơ) là \sdGenGV{} điểm.
Vì vậy mỗi biến thể được sinh ba lượt, mỗi lượt chấm một lần; một câu được tính là đúng khi phiếu \emph{accurate} chiếm đa số tuyệt đối có trọng số.
Riêng MiniRAG tái lập chỉ có \emph{một} lượt sinh, và hệ quả của điều này được nêu ở mỗi phép so liên quan.
